Визуализация течения жидкости в компьютерной графике опирается на различные
численные и графические методы. Для систематического анализа удобно ввести
классификацию алгоритмов по двум основным основаниям:
\begin{itemize}
  \item по \textbf{представлению среды} (эйлеровы, лагранжевы и гибридные);
  \item по \textbf{степени физической достоверности} (физически обоснованные и эвристические методы).
\end{itemize}

\subsection{Классификация по способу представления среды}

С точки зрения математического описания и хранения данных о состоянии
жидкости в памяти компьютера, выделяют три основных подхода \cite{chinapaper}.

\subsubsection{Эйлеровы (сеточные) методы}

В эйлеровом подходе жидкость рассматривается как континуальная
среда, а её состояние описывается полями величин на фиксированной пространственной
сетке. На узлах или в ячейках сетки хранятся компоненты вектора скорости
$u = (u_x, u_y, u_z)$, давление $p$, плотность $\rho$ и другие
параметры.

Численное решение уравнений движения жидкости в эйлеровом подходе
сводится к последовательному обновлению значений этих полей по шагам времени.

\subsubsection{Лагранжевы (частичные) методы}

В лагранжевом подходе вместо фиксированной сетки рассматривается
набор дискретных частиц, которые движутся одним потоком. Каждая частица
характеризуется положением $x_i$, скоростью $v_i$, массой
$m_i$, а также, при необходимости, плотностью, давлением и другими атрибутами.

Движение частиц задаётся интегрированием уравнений движения во времени.
Частицы могут интерпретироваться как материальные точки или как носители
некоторого объёма жидкости. Подобные методы широко применяются, например,
в алгоритмах сглаженной гидродинамики частиц (Smoothed Particle Hydrodynamics, SPH)
и позиционно-ориентированной динамики (Position-Based Fluids, PBF).

\subsubsection{Гибридные методы}

Гибридные методы сочетают эйлеровы и лагранжевы подходы.
Типичный пример \textemdash{} методы типа PIC/FLIP (Particle-in-Cell / Fluid-Implicit-Particle),
в которых:
\begin{itemize}
  \item сетка используется для решения глобальных задач (например, для обеспечения
        несжимаемости и вычисления давления);
  \item частицы используются для хранения массы и импульса, а также для
        повышения детализации (брызги, капли, мелкие вихри).
\end{itemize}

Такая комбинация позволяет одновременно воспользоваться устойчивостью и
простотой сеточных схем и гибкостью частичных методов при описании сложных
поверхностей и локальных эффектов.

\subsection{Классификация по уровню физической достоверности}

Второй важный аспект классификации \textemdash{} это степень, в которой алгоритм
соответствует физическим законам течения жидкости.

\subsubsection{Физически обоснованные методы}

Физически обоснованные методы основываются на уравнениях гидродинамики
и стремятся к точному или приближенному решению этих уравнений.

\paragraph{Численное решение уравнений Навье-Стокса} является классическим
подходом к моделированию течения несжимаемой жидкости \cite{holybook}.

Для вязкой несжимаемой жидкости
(\(\rho = \mathrm{const}\)) систему можно записать как:
\begin{equation}
  \begin{cases}
    \displaystyle
    \frac{\partial \mathbf{u}}{\partial t}
    =
    - (\mathbf{u}\cdot\nabla)\mathbf{u}
    -\frac{1}{\rho}\,\nabla p
    + \nu\nabla^{2}\mathbf{u}
    + \mathbf{f}, \\[6pt]
    \displaystyle
    \nabla\cdot\mathbf{u} = 0.
  \end{cases}
  \label{eq:NS}
\end{equation}

Здесь
\begin{itemize}
  \item \(\mathbf{u}(\mathbf{x},t)\) \textemdash{} поле скорости жидкости,
        \(u=(u_x,u_y,u_z)\);
  \item \(p(\mathbf{x},t)\) \textemdash{} поле давления;
  \item \(\rho\) \textemdash{} плотность (для несжимаемой жидкости считаем
        постоянной);
  \item \(\nu\) \textemdash{} кинематическая вязкость (\(\nu=\eta/\rho\), где \(\eta\) \textemdash{} коэффициент динамической вязкости);
  \item \(\mathbf{f}(\mathbf{x},t)\) \textemdash{} вектор внешних сил
        на единицу массы (например, гравитация \(\mathbf{g}\)).
\end{itemize}

Оператор набла \(\nabla\) в декартовых координатах имеет вид
\begin{equation}
  \nabla =
  \left(
  \frac{\partial}{\partial x},
  \frac{\partial}{\partial y},
  \frac{\partial}{\partial z}
  \right).
\end{equation}
В зависимости от типа поля он задаёт:
\begin{itemize}
  \item \textit{градиент} скалярного поля \(\phi(x,t)\):
        \begin{equation}
          \nabla\phi =
          \left(
          \frac{\partial\phi}{\partial x},
          \frac{\partial\phi}{\partial y},
          \frac{\partial\phi}{\partial z}
          \right);
        \end{equation}
  \item \textit{дивергенцию} векторного поля \(u\):
        \begin{equation}
          \nabla\cdot u
          =
          \frac{\partial u_x}{\partial x}
          + \frac{\partial u_y}{\partial y}
          + \frac{\partial u_z}{\partial z};
        \end{equation}
  \item \textit{лапласиан} (оператор диффузии)
        скалярного или векторного поля:
        \begin{equation}
          \nabla^{2}\phi
          =
          \frac{\partial^{2}\phi}{\partial x^{2}}
          +
          \frac{\partial^{2}\phi}{\partial y^{2}}
          +
          \frac{\partial^{2}\phi}{\partial z^{2}},
        \end{equation}
        а для векторного поля \(u\) лапласиан применяется
        покомпонентно:
        \(\nabla^{2}u
        = (\nabla^{2}u_x,\nabla^{2}u_y,\nabla^{2}u_z)\).
\end{itemize}

Условие \(\nabla\cdot u=0\) в~\eqref{eq:NS} называется уравнением
непрерывности и
означает несжимаемость: локально объём жидкости не меняется,
масса не создаётся и не исчезает.


При численном решении системы \eqref{eq:NS} ключевую роль играет
уравнение Пуассона для давления\footnote{Подробнее про уравнение Пуассона \url{https://en.wikipedia.org/wiki/Poisson\%27s\_equation\#Fluid\_dynamics}}. Его получают, применяя
оператор дивергенции к уравнению движения и используя условие
несжимаемости \(\nabla\cdot u=0\). В результате для поля
давления \(p(x,t)\) получается эллиптическое уравнение
вида:
\begin{equation}
  \nabla^{2} p
  =
  \rho\,
  \nabla\cdot\left(
  -(u\cdot\nabla)u
  + f
  \right),
  \label{eq:poisson-pressure}
\end{equation}
где правая часть задаётся известными на текущем шаге скоростями
и внешними силами.

Уравнение \eqref{eq:poisson-pressure} определяет такое
распределение давления, которое подправляет поле скорости так,
чтобы после шага интегрирования по времени выполнялось условие
несжимаемости. Во многих численных схемах сначала вычисляют промежуточную скорость без учёта давления,
затем решают уравнение Пуассона \eqref{eq:poisson-pressure} и,
используя найденное поле \(p\), проецируют скорость на
соленоидальное подпространство.

\paragraph{Метод решёточных уравнений Больцмана (LBM)} является альтернативным подходом к моделированию, отличающимся от прямого численного
решения уравнений Навье-Стокса. В отличие от традиционных методов вычислительной гидродинамики,
работающих на макроскопическом уровне, LBM оперирует на мезоскопическом уровне,
основываясь на кинетической теории газов.

Вместо непосредственного вычисления полей давления и скорости, метод
моделирует эволюцию функции распределения частиц \(f_i(x, t)\)
на дискретной решётке. Основное уравнение LBM с оператором столкновений
(часто используется приближение Бхатнагара-Гросса-Крука)
выглядит следующим образом:
\begin{equation}
  f_i(x + c_i \Delta t,\, t + \Delta t)
  = f_i(x, t)
  - \frac{\Delta t}{\tau}
  \bigl(f_i(x, t) - f_i^{eq}(x, t)\bigr),
  \label{eq:LBM}
\end{equation}
где:
\begin{itemize}
  \item \(f_i(x, t)\) \textemdash{} функция распределения частиц,
        движущихся в направлении \(i\);
  \item \(c_i\) \textemdash{} дискретный набор векторов скоростей,
        определяемый топологией решётки (например, D2Q9 или D3Q19);
  \item \(\tau\) \textemdash{} время релаксации, связанное с кинематической
        вязкостью жидкости;
  \item \(f_i^{eq}(x, t)\) \textemdash{} локальная равновесная
        функция распределения.
\end{itemize}

Макроскопические величины восстанавливаются как моменты функций
распределения:
\[
  \rho(x, t)
  = \sum_i f_i(x, t),
  \qquad
  \rho(x, t)\, u(x, t)
  = \sum_i f_i(x, t)\, c_i .
\]
Показано, что в гидродинамическом пределе (при малых числах Кнудсена
и Маха) из уравнения~\eqref{eq:LBM} с помощью разложения
Чепмена\textemdash{}Энского восстанавливаются уравнения Навье\textemdash{}Стокса.
Благодаря полностью локальному характеру вычислений метод особенно
эффективен при реализации на графических процессорах (GPU).
\bigskip

Визуализация в физических методах является прямым следствием численного решения
этих уравнений (или их упрощённых форм). К ним относятся:
эйлеровские схемы на сетке (включая метод stable fluids), SPH,
гибридные подходы PIC/FLIP и др. Основное достоинство таких методов \textemdash{}
высокая физическая достоверность и предсказуемое поведение модели при
изменении параметров.

\subsubsection{Эвристические и художественные методы}

\textbf{Эвристические} (или художественные) методы не стремятся к строгому
решению уравнений движения жидкости. Вместо этого они используют различные
визуальные приёмы, которые создают у наблюдателя впечатление течения
жидкости. К таким приёмам можно отнести:
\begin{itemize}
  \item процедурные текстуры и шумы (например, шум Перлина, curl-noise)
        для имитации турбулентности;
  \item экранно-пространственные эффекты (screen-space advection), при
        которых "течет" уже отрендерованное изображение;
  \item простые модели адвекции и диффузии для скалярных полей без строгого
        учёта несжимаемости.
\end{itemize}

Подобные алгоритмы, как правило, проще в реализации и дешевле по
вычислительным ресурсам, что делает их популярными в реальном времени
(игровые движки, веб-визуализация). Однако они не гарантируют физической
корректности и предназначены прежде всего для получения визуально
